
\section{Rectificador de media onda}

\subsection{Instrumentos y componentes utilizados}

\begin{itemize}
    \item Fuente de tensión alterna de $12\ \mathrm{V_{RMS}}$, $V_S$, y
    $50\ \mathrm{Hz}$, $f_S$.
    \item Osciloscopio.
    \item Multímetro digital.
    \item Diodo rectificador, $D$.
    \item Fusible de protección conectado en serie con la entrada, $F$.
    \item Resistencia limitadora de corriente de valor nominal
    $R_{\lim}=4{,}7\ \Omega$.
    \item Resistencia de carga de valor nominal
    $R_L=100\ \Omega$.
    \item Tres capacitores electrolíticos: $C_1 =100\ \mu\mathrm{F}$, $C_2 =220\ \mu\mathrm{F}$, $C_3 =2200\ \mu\mathrm{F}$.
    \item Cables para el conexionado.
\end{itemize}


\subsection{Objetivos}

El objetivo de esta experiencia fue estudiar experimentalmente el
funcionamiento de un circuito rectificador de media onda.

En particular, se buscó observar la señal de salida sobre la resistencia obtenida a partir
de una entrada sinusoidal,$V_S$.  Luego utilizando el osciloscopio y el multímetro determinar las caracteristicas del circuito y estudiar la modificación de la tensión de salida al
incorporar capacitores de filtrado de diferentes capacidades.

Además, se midió la amplitud pico a pico del rizado para cada capacitor
y se calculó analíticamente el factor de rizado correspondiente a la
señal de salida sin capacitor.

\subsection{Desarrollo experimental}

Para esta experiencia se implementó un circuito rectificador de media
onda. Como medidas de protección, se conectaron antes del diodo un fusible y
una resistencia limitadora de corriente $R_{\lim}$. El fusible permite
interrumpir el circuito ante una corriente excesiva, mientras que la
resistencia $R_{\lim}$ limita la corriente máxima que puede circular por
el circuito.

Inicialmente se armó el circuito rectificador de media onda sin conectar
el capacitor de filtrado como se ve en el esquemático de la Figura \ref{fig:rectificador_media_onda}.

\begin{figure}[H]
    \centering
    \begin{tikzpicture}
        % Paths, nodes and wires:
        \draw (4.5, 3.75) to[american resistor, l={$R_{lim}$}] (2, 3.75);
        \draw (4, 3.75) to[empty diode, l={$D$}] (6.5, 3.75);
        \draw (6.5, 3.75) to[american resistor, l={$R_{L}$}] (6.5, 1);

        \draw (0.5, 3.75) to[sinusoidal current source, l={$V_S$}] (0.5, 1);
        \draw (6.5, 1) -- (0.5, 1) -- (2, 1);
        \draw (0.5, 3.75) to[generic, l={$F$}] (2, 3.75);
    \end{tikzpicture}
    \caption{Circuito rectificador de media onda sin efecto capacitivo.}
    \label{fig:rectificador_media_onda}
\end{figure}

Luego se conectaron en paralelo con la resistencia de carga
los capacitores de a uno por vez, como se puede ver en el esquemático de la Figura \ref{fig:rectificador_media_onda_cap}.

\begin{figure}[H]
    \centering
    \begin{tikzpicture}
        % Paths, nodes and wires:
        \draw (0.5, 3.75) to[generic, l={$F$}] (2, 3.75);
        \draw (2, 3.75) to[american resistor, l={$R_{lim}$}] (4.5, 3.75);
        \draw (4.5, 3.75) to[empty diode, l={$D$}] (7, 3.75);
        % Nodo de unión arriba
        \draw (7, 3.75) -- (8.5, 3.75);
        % R_L con voltaje
        \draw (7, 3.75) to[american resistor, l={$R_{L}$}] (7, 1);
        % Capacitor en paralelo con R_L
        \draw (8.5, 3.75) to[C, l={$C_i$}] (8.5, 1);
        % Nodo de unión abajo
        \draw (8.5, 1) -- (7, 1);
        % Fuente de tensión
        \draw (0.5, 3.75) to[sinusoidal voltage source, l={$V_S$}] (0.5, 1);
        % Cable de retorno (masa)
        \draw (0.5, 1) -- (8.5, 1);
    \end{tikzpicture}
    \caption{Circuito rectificador de media onda con el capacitor.}
    \label{fig:rectificador_media_onda_cap}
\end{figure}

Para cada configuración Se conectó un canal del osciloscopio a la salida, sobre la resistencia de carga. De esta manera se
observo la señal
rectificada de salida. Ademas, se registró la amplitud pico a pico del rizado,
el valor medio indicado por el osciloscopio y la tensión continua medida
con el multímetro.

Al utilizar capacitores electrolíticos se respetó su polaridad,
conectando el terminal positivo al nodo positivo de la salida y el
terminal negativo al nodo de referencia del circuito.

\subsection{Datos obtenidos}

Antes de realizar las mediciones de la tensión de salida se verificaron
los valores de la fuente y de las resistencias utilizadas.

Los valores nominales y medidos se presentan en el Cuadro
\ref{tab:parametros_media_onda}.

\begin{table}[H]
    \centering
    \renewcommand{\arraystretch}{1.4}
    \setlength{\arrayrulewidth}{0.5pt}

    \begin{tabular}{|
        >{\centering\arraybackslash}m{4.5cm}|
        >{\centering\arraybackslash}m{3cm}|
        >{\centering\arraybackslash}m{3cm}|}
        \hline

        \rowcolor{headercolor}
        \textcolor{white}{\textbf{Magnitud}} &
        \textcolor{white}{\textbf{Valor nominal}} &
        \textcolor{white}{\textbf{Valor medido}} \\
        \hline

        \rowcolor{rowcolor}
        $R_L \ [\Omega]$ &
        $100$ &
        $104$ \\
        \hline

        $V_S \ [V rms]$ &
        $12$ &
        $13$ \\
        \hline

        \rowcolor{rowcolor}
         $f_S \ [\mathrm{Hz}]$ &
        $50$ &
        -- \\
        \hline

        $R_{\lim}$ &
        $4{,}7\ \Omega$ &
        $7{,}4\ \Omega$ \\
        \hline

    \end{tabular}

    \caption{Valores nominales y medidos de los componentes utilizados
    en el rectificador de media onda.}
    \label{tab:parametros_media_onda}
\end{table}

Las mediciones realizadas sobre la salida del rectificador se presentan
en el Cuadro \ref{tab:mediciones_media_onda}.

\begin{table}[H]
    \centering
    \renewcommand{\arraystretch}{1.4}
    \setlength{\arrayrulewidth}{0.5pt}

    \resizebox{\columnwidth}{!}{%
    \begin{tabular}{|
        >{\centering\arraybackslash}m{3.3cm}|
        >{\centering\arraybackslash}m{2.5cm}|
        >{\centering\arraybackslash}m{2.8cm}|
        >{\centering\arraybackslash}m{3cm}|
        >{\centering\arraybackslash}m{3.3cm}|}
        \hline

        \rowcolor{headercolor}
        \textcolor{white}{\textbf{Configuración}} &
        \textcolor{white}{\textbf{$C$ [$\mu$F]}} &
        \textcolor{white}{\textbf{$V_{pp}$ [V]}} &
        \textcolor{white}{\textbf{$V_{\mathrm{medio}}$ [V]}} &
        \textcolor{white}{\textbf{$V_{\mathrm{DC}}$ [V]}} \\
        \hline

        \rowcolor{rowcolor}
        Sin capacitor &
        -- &
        16,0 &
        4,85 &
        5,07 \\
        \hline

        Capacitor 1 &
        100 &
        13,2 &
        8,50 &
        8,75 \\
        \hline

        \rowcolor{rowcolor}
        Capacitor 2 &
        220 &
        9,0 &
        10,80 &
        11,00 \\
        \hline

        Capacitor 3 &
        2200 &
        1,5 &
        12,60 &
        12,90 \\
        \hline

    \end{tabular}%
    }

    \caption{Mediciones de la tensión de salida del rectificador de
    media onda.}
    \label{tab:mediciones_media_onda}
\end{table}

En el Cuadro \ref{tab:mediciones_media_onda}, $V_{pp}$ corresponde a la
amplitud pico a pico de la señal de salida o del rizado,
$V_{\mathrm{medio}}$ corresponde al valor medio indicado por el
osciloscopio y $V_{\mathrm{DC}}$ corresponde a la tensión medida con el
multímetro en la escala de continua.

\subsection{Analisis de la salida del circuito sin capacitor}

\subsubsection{Amplitud cero a pico}

En el rectificador de media onda sin capacitor, el valor mínimo de la
señal de salida es aproximadamente igual a cero. Por lo tanto, la
amplitud pico a pico coincide aproximadamente con la amplitud cero a
pico:

\begin{equation}
    V_{op}\approx V_{pp}.
\end{equation}

A partir de la medición realizada con el osciloscopio se obtuvo:

\begin{equation}
    V_{op}=V_{pp}=16{,}0\ \mathrm{V}.
\end{equation}

Por otro lado, considerando el valor eficaz medido de la fuente,
$V_S=13\ \mathrm{V_{RMS}}$, su amplitud pico ideal es:

\begin{align*}
    V_{S,p}
    &=V_S\sqrt{2}\ =18{,}38\ \mathrm{V}.
\end{align*}

La amplitud de salida resulta menor que el pico de la fuente debido a
la caída de tensión en el diodo y en la resistencia limitadora
$R_{\min}$.

\subsubsection{Valor eficaz de la tensión de salida}

La tensión de salida del rectificador de media onda sin capacitor puede
expresarse como:

\begin{equation}
    v_o(t)=
    \begin{cases}
        V_{op}\sin(\omega t),
        & 0\leq t\leq \dfrac{T}{2},\\[6pt]
        0,
        & \dfrac{T}{2}<t<T.
    \end{cases}
\end{equation}

El valor eficaz de la señal de salida se calcula mediante:

\begin{equation}
    V_{\mathrm{RMS}}
    =
    \sqrt{
        \frac{1}{T}
        \int_0^{T/2}
        V_{op}^{\,2}\sin^2(\omega t)\,dt
    }.
\end{equation}

De esta manera:

\begin{align*}
    V_{\mathrm{RMS}}
    &=
    \frac{V_{op}}{2}.
\end{align*}

Por lo tanto:

\begin{equation}
    \boxed{V_{\mathrm{RMS}}=8{,}0\ \mathrm{V}}.
\end{equation}

El valor calculado como $V_{RMS}$ no reprecenta lo mismo que el valor registrado mediante el multímetro, por lo que da valores significativamente diferentes. EL voltaje con el multímetro fue medido en la escala de
tensión continua. En consecuencia, este valor representa la componente
continua o valor medio de la señal de salida, y no el valor eficaz total
de la señal.

\subsubsection{Valor medio de la tensión de salida}

El valor medio teórico de la señal rectificada de media onda se calcula
como:

\begin{equation}
    V_{\mathrm{medio}}
    =
    \frac{1}{T}
    \int_0^{T/2}
    V_{op}\sin(\omega t)\,dt.
\end{equation}

se obtiene:

\begin{align*}
    V_{\mathrm{medio}}
    &=
    \frac{V_{op}}{\pi}.
\end{align*}

Reemplazando el valor de la amplitud cero a pico:

\begin{equation}
    \boxed{V_{\mathrm{medio,teo}}=5{,}09\ \mathrm{V}}.
\end{equation}

Por otro lado el valor medido mediante el multímetro en la escala de continua fue de $5{,}07 \ \mathrm{V}$ y el valor medio registrado mediante el osciloscopio fue de $4{,}85\ \mathrm{V}$. La diferencia porcentual entre el valor medio teórico y el valor medido practico en cada caso se calculó mediante al ecuación:

\begin{equation}
    \varepsilon_{\mathrm{mult}}
    =
    \frac{
        \left|
        V_{\mathrm{DC,med}}-V_{\mathrm{medio,teo}}
        \right|
    }{
        V_{\mathrm{medio,teo}}
    }
    \cdot100.
\end{equation}

La diferencia porcentual entre el valor medio teorico y el medido con el multimetro es de $0{,}45\%$, Esta pequeña discrepancia puede asociarse a la precisión de los instrumentos, a las tolerancias de los componentes y a que el modelo teórico considera un diodo y una fuente ideales.

Por otra parte, la diferencia entre el valor teórico y el valor medio indicado por el osciloscopio fue de $4{,}72\%$. Esta diferencia no debe atribuirse exclusivamente a la precisión del osciloscopio, sino también a que la señal experimental no es una media sinusoide ideal. La tensión umbral del diodo, la resistencia limitadora, la impedancia interna de la fuente y la resistencia de carga pueden modificar la forma de onda.


\subsubsection{Factor de rizado}

La señal de salida, $v_o(t)$, puede interpretarse como la suma de una
componente continua y una componente variable o de rizado:

\begin{equation}
    v_o(t)
    =
    V_{\mathrm{medio}}
    +
    v_{\mathrm{rizado}}(t).
\end{equation}

El valor eficaz total de la señal se relaciona con su componente
continua y con el valor eficaz de la componente de rizado mediante:

\begin{equation}
    V_{\mathrm{RMS}}^2
    =
    V_{\mathrm{medio}}^2
    +
    V_{\mathrm{rizado,RMS}}^2.
\end{equation}

Por lo tanto, el valor eficaz de la componente de rizado puede
calcularse como:

\begin{equation}
    V_{\mathrm{rizado,RMS}}
    =
    \sqrt{
        V_{\mathrm{RMS}}^2
        -
        V_{\mathrm{medio}}^2
    }.
\end{equation}

El factor de rizado se define como la relación entre el valor eficaz de
la componente alterna y el valor de la componente continua:

\begin{equation}
    r
    =
    \frac{
        V_{\mathrm{rizado,RMS}}
    }{
        V_{\mathrm{medio}}
    }.
\end{equation}

Utilizando el valor medio registrado por el osciloscopio,
 ,$4{,}85\ \mathrm{V}$, y el valor eficaz
calculado previamente, $8{,}0\ \mathrm{V}$, se obtiene $V_{\mathrm{rizado,RMS,osc}} =6{,}36\ \mathrm{V}$. En consecuencia, el factor de rizado estimado a partir de la medición
del osciloscopio,$r_{\mathrm{osc}}$, resulta  $1{,}31$. 

Por otro lado, utilizando el valor medio medido con el multímetro
digital, $5{,}07\ \mathrm{V}$, se obtiene $V_{\mathrm{rizado,RMS,mult}}=6{,}18\ \mathrm{V}$. 
Por lo tanto $r_{\mathrm{mult}}=1{,}22$ .

En ambos casos el factor de rizado es mayor que uno. Esto indica que el
valor eficaz de la componente alterna o de rizado es mayor que el valor
de la componente continua de la señal de salida. Por lo tanto, la salida
del rectificador de media onda sin capacitor presenta una componente
variable considerable y se encuentra alejada de una tensión continua
constante.

La proporción de la potencia de salida asociada a la componente
continua puede estimarse mediante:

\begin{equation}
    \eta
    =
    \frac{P_{\mathrm{DC}}}{P_{\mathrm{total}}}
    =
    \frac{V_{\mathrm{medio}}^2}
    {V_{\mathrm{RMS}}^2}.
\end{equation}

Utilizando el valor medio medido con el multímetro $\eta_{\mathrm{mult}}=0{,}402$ . Por lo tanto, aproximadamente el $40{,}2\%$ de la potencia total de salida se encuentra asociada a la componente continua, mientras que el $59{,}8\%$ restante corresponde a la componente de rizado. 
Utilizando el valor medio indicado por el osciloscopio $\eta_{\mathrm{osc}}=0{,}368$. En este caso, la proporción estimada de potencia continua resulta del
$36{,}8\%$.

La diferencia entre ambos valores no indica que el circuito haya tenido
pérdidas distintas durante cada medición, ya que ambos instrumentos
analizaron la misma señal. La discrepancia surge de la diferencia entre
los valores medios registrados por cada instrumento y de las
condiciones particulares de cada método de medición.

Por lo tanto, menos del $40{,}5\%$ de la potencia de salida de un rectificador de media onda sin filtrado corresponde a su componente continua. Esto implica que, si el objetivo es alimentar una carga exclusivamente con corriente continua, más de la mitad de la potencia de salida se encuentra asociada a la componente alterna o de rizado y no resulta útil para dicho propósito. 


\subsection{Analisis de la salida del circuito con capacitor}
Hacer el inciso 2, 3, 4, 5

\newpage

